\subsection{Halluzinationen}
Als Halluzinationen bezeichnet man im kontext von Large Language Models (LLMs) die generierung von Informationen, welche nicht auf den Trainingsdaten basieren. 
Diese Informationen können sprachlich plausibel und korrekt erscheinen, sind jedoch faktisch nicht durch die verwendeten Trainingsdaten gedeckt. 
Dieses Phänomen wird im bereich der Natural Language Generation (NLG) als Halluzination bezeichnet.\autocite[3]{jiSurveyHallucinationNatural2023}
\newline \noindent
Desweiteren unterscheidet man zwischen zwei Arten von Halluzinationen:
\begin{enumerate}
    \item \textbf{Intrinsic Halluzinationen:} Hier kann der generierte Inhalt teilweise durch die quelle gestützt werden. \autocite[3]{jiSurveyHallucinationNatural2023}
    \item \textbf{Extrinsic Halluzinationen:} Hier kann der generierte Inhalt nicht durch die Quelle gestützt werden. \autocite[3]{jiSurveyHallucinationNatural2023}
\end{enumerate}

\subsubsection{Terminologie}
%%%On faithfulness and factuality in abstractive summarization. In Proceedings of the 58th Annual Meeting of the Association for Computational Linguistics
\begin{enumerate}
    \item Faithfulness: Informationen konsistent und korrekt aus den gegebenen Quellen wiedergeben. \autocite[5]{jiSurveyHallucinationNatural2023}
    \item Factuality: Beschreibt die Faktische korrektheit des generierten Inhalts, basierend auf einem gegebenen Wissensstand. 
            Es beschreibt somit auch eine Faktenbasierte generierung der Antwort. Bei einem Factuality ist der fact als art 'Welt Wissen' zu verstehen. \autocite[5]{jiSurveyHallucinationNatural2023}

        
        \end{enumerate}
\subsubsection{Cause}

Note:
- Relevanz nicht nur für nutzer sondern systeme die llms einsetzen und vollautomatisch laufen. Halluzinierte antworten können zu großen problemen führen.